\section{Discussion}
\label{sec:discussion}

\subsection{Yogac\=ara as Architectural Guide}

The companion survey~\cite{hui2026taxonomy} identifies four cases where
Yogac\=ara concepts predict design choices that the AI memory community has
independently converged upon. \alaya{} was designed with these predictions
as explicit architectural constraints rather than post-hoc rationalizations:

\begin{itemize}[nosep]
    \item \textbf{Heterogeneous ripening} (\emph{vip\=aka}) predicted that
    consolidation should produce outputs qualitatively different from inputs.
    \alaya{}'s consolidation converts time-bound episodes into
    decontextualized, typed semantic nodes---a structural transformation,
    not a summary.
    \item \textbf{Gradual perfuming} (\emph{v\=asan\=a}) predicted that
    preferences should emerge from accumulation, not extraction. \alaya{}'s
    five-impression threshold before preference crystallization implements
    this directly.
    \item \textbf{Transformation of the basis}
    (\emph{\=a\'sraya-par\=av\d{r}tti}) predicted that the memory store
    requires periodic reorganization. \alaya{}'s \texttt{transform()}
    operation---deduplication, contradiction resolution, category
    discovery---implements this prediction.
    \item \textbf{Seed mutual conditioning} predicted that stored elements
    should activate related elements through weighted, use-adaptive links.
    \alaya{}'s Hebbian graph overlay with spreading activation implements
    this structural isomorphism directly.
\end{itemize}

The value of Yogac\=ara as an architectural guide is not mystical but
structural: it provides a coherent theory of memory-as-process that
identifies operations (perfuming, ripening, transformation) which neither
pure neuroscience models nor pure engineering intuition would prioritize
independently. The preference emergence mechanism, in particular, was
designed from the v\=asan\=a concept before we found empirical support in
the memory systems literature.

\subsection{Graceful Degradation}

A deliberate design principle is that \alaya{} never fails due to missing
optional capabilities. The system operates in three degradation modes:

\begin{enumerate}[nosep]
    \item \textbf{Full capability:} Embeddings + LLM provider available.
    All retrieval signals (BM25, vector, graph), embedding-based clustering,
    LLM-generated category labels.
    \item \textbf{No embeddings:} BM25 + graph retrieval only. Category
    discovery uses graph-only clustering (shared Hebbian link overlap).
    \item \textbf{No LLM:} Consolidation is unavailable, but all other
    lifecycle operations function: forgetting, transformation (on existing
    semantic nodes), category discovery (from existing nodes), preference
    emergence (from accumulated impressions).
\end{enumerate}

This design reflects a practical reality: many deployment contexts
(offline agents, privacy-sensitive environments, cost-constrained settings)
cannot guarantee access to embedding models or LLMs for every memory
operation. The system should do the best it can with available resources
rather than failing entirely.

\subsection{Positioning in the Landscape}

The companion survey's capability scoring~\cite{hui2026taxonomy} evaluates
systems across eight discriminating features. \alaya{} implements all eight:
episodic memory, semantic memory, graph structure, forgetting,
consolidation, contradiction resolution, preference learning, and hybrid
retrieval fusion. Among the 88 surveyed systems, only 2.3\% score 8/8 on
this metric.

The distinctive characteristics relative to comparable systems are:
\begin{itemize}[nosep]
    \item vs.\ \textbf{Mem0}~\cite{choudhary2025mem0}: \alaya{} adds
    principled forgetting, graph-based retrieval, emergent (not extracted)
    preferences, and consolidation. Mem0 requires external services
    (Qdrant/Neo4j); \alaya{} requires none.
    \item vs.\ \textbf{Zep/Graphiti}~\cite{dempsey2025graphiti}: Both use
    knowledge graphs, but \alaya{}'s graph is Hebbian (reshapes through use)
    while Graphiti's is LLM-extracted (static after creation). \alaya{} adds
    forgetting and preference emergence.
    \item vs.\ \textbf{SYNAPSE}~\cite{jiang2026synapse}: Both implement
    CLS-inspired consolidation and activation-based retrieval. \alaya{} adds
    emergent ontology and the implicit (preference) store. SYNAPSE is
    evaluated on specific benchmarks; \alaya{} is a general-purpose library.
    \item vs.\ \textbf{Generative Agents}~\cite{park2023generative}: Both
    implement episodic-to-semantic consolidation (reflection). \alaya{} adds
    principled forgetting (Bjork vs.\ recency weighting), Hebbian graph
    evolution, and emergent categories.
\end{itemize}

\subsection{Limitations}

\alaya{} v0.1.0 has several limitations:
\begin{itemize}[nosep]
    \item \textbf{No end-to-end benchmark results:} We report baseline
    crossover analysis but have not yet evaluated \alaya{}'s full pipeline
    against LoCoMo/LongMemEval. This is ongoing work.
    \item \textbf{Flat categories only:} The emergent ontology discovers
    flat categories without hierarchy. Category hierarchy, prototype theory,
    and category seeds are planned for v0.2.0.
    \item \textbf{Single-agent only:} No multi-agent memory sharing,
    ownership semantics, or federated memory.
    \item \textbf{Thresholds are hand-tuned:} Cosine similarity thresholds
    for category assignment (0.6), discovery (0.7), and merge (0.85) are
    based on reasonable defaults, not learned from data.
    \item \textbf{SQLite scalability:} While SQLite handles millions of
    records efficiently, the in-process architecture limits concurrent write
    access to a single process.
\end{itemize}
